% from ega4-17-fr.tex
\oldpage[IV-4]{75}

close de $k$, $X\otimes_k k'$ est alors lisse sur $k'$, et en vertu de (17.7.3, (ii)), il suffit de prouver que $Y\otimes_k k'$ est lisse sur $k'$ (puisque $\Spec(k')\to\Spec(k)$ est fidèlement plat); on peut donc se borner au cas où $k$ est \emph{algébriquement clos}. Comme l'ensemble des points de $Y$ rationnels sur $k$ est alors très dense dans $Y$ (10.4.8) et comme l'ensemble des points de $Y$ où $Y$ est lisse sur $k$ est ouvert, on voit qu'il suffit de prouver que $Y$ est lisse sur $k$ en tout point rationnel sur $k$. Mais en un tel point $y$, dire que $Y$ est lisse sur $k$ en ce point équivaut à dire que $Y$ est \emph{régulier} en $y$ (17.5.1 et 6.7.8). Or on a $y=g(x)$ pour un $x\in X$ et par hypothèse (17.5.1) $X$ est régulier au point $x$; comme $X$ et $Y$ sont alors localement noethériens et que $g$ est plat, $Y$ est bien régulier au point $y$ (6.5.1, (i)).

(ii) \emph{Cas des morphismes étales}. Par (i) on sait déjà que $f$ est lisse au point $y$; en vertu de (17.6.1), il suffit donc de montrer que $f$ est quasi-fini au point $y$, ou encore que $\sh{O}_{Y,y}$ est une $k$-algèbre finie. Or, comme $\sh{O}_{X,x}$ est un $\sh{O}_{Y,y}$-module fidèlement plat ($0_{\mathrm I}$, 6.6.2), $\sh{O}_{Y,y}$ s'identifie à un sous-$k$-module de $\sh{O}_{X,x}$ et comme $\sh{O}_{X,x}$ est par hypothèse une $k$-algèbre finie, il en est de même de $\sh{O}_{Y,y}$.

\begin{proposition}[17.7.8]
\label{IV.17.7.8}
Les notations étant celles de (8.8.1), on suppose $X_\alpha$ et $Y_\alpha$ localement de présentation finie sur $S_\alpha$. Soient $f_\alpha:X_\alpha\to Y_\alpha$ un $S_\alpha$-morphisme, $f:X\to Y$ le $S$-morphisme correspondant.
\begin{enumerate}
  \item[(i)] Soient $x$ un point de $X$, $x_\lambda$ sa projection canonique dans $X_\lambda$. Pour que $f$ soit lisse (resp. non ramifié, resp. étale) au point $x$, il faut et il suffit qu'il existe $\lambda\geq\alpha$ tel que $f_\lambda$ soit lisse (resp. non ramifié, resp. étale) au point $x_\lambda$.
  \item[(ii)] Supposons de plus $X_\alpha$ quasi-compact. Pour que $f$ soit lisse (resp. non ramifié, resp. étale), il faut et il suffit qu'il existe $\lambda\geq\alpha$ tel que $f_\lambda$ soit lisse (resp. non ramifié, resp. étale).
\end{enumerate}
\end{proposition}

(i) Si $y=f(x)$, $y_\lambda=f_\lambda(x_\lambda)$ est la projection canonique de $y$ dans $Y_\lambda$, et l'on a $f^{-1}(y)=f_\lambda^{-1}(y_\lambda)\otimes_{k(y_\lambda)}k(y)$; la partie de l'énoncé concernant les morphismes non ramifiés résulte donc de (17.7.1, (i)), et il suffit donc de considérer le cas des morphismes lisses. Comme $f$ et $f_\lambda$ sont localement de présentation finie, il revient au même de dire que $f^{-1}(y)$ est géométriquement régulier au point $x$ ou que $f_\lambda^{-1}(y_\lambda)$ est géométriquement régulier au point $x_\lambda$ (6.7.8). La proposition résulte donc de (17.5.1) et de (11.2.6).

(ii) Pour tout $\lambda$, soit $U_\lambda$ l'ensemble ouvert des $x_\lambda\in X_\lambda$ tels que $f_\lambda$ soit lisse (resp. non ramifié, resp. étale) au point $x_\lambda$; soit $V_\lambda$ son image réciproque dans $X$. Comme, par hypothèse, pour tout $x\in X$ il existe, en vertu de (i), un $\lambda$ tel que $f_\lambda$ soit lisse (resp. non ramifié, resp. étale) au point $x_\lambda$, $X$ est réunion des $V_\lambda$. D'ailleurs (17.3.3), pour $\lambda\leq\mu$ on a $V_\lambda\subset V_\mu$, donc, comme $X$ est quasi-compact, il existe un indice $\mu$ tel que $X=V_\mu$. Comme les $X_\lambda$ sont quasi-compacts, il résulte alors de (8.3.4) qu'il existe un indice $\nu\geq\mu$ tel que l'image réciproque de $U_\mu$ dans $X_\nu$ soit $X_\nu$ tout entier, ce qui signifie que $f_\nu$ est lisse (resp. non ramifié, resp. étale) par (17.3.3).

\begin{corollary}[17.7.9]
\label{IV.17.7.9}
Soient $S=\Spec(A)$ un schéma affine, $f:X\to S$ un morphisme. Les conditions suivantes sont équivalentes:
\begin{enumerate}
  \item[a)] $f$ est un morphisme de présentation finie et lisse (resp. non ramifié, resp. étale).

  